\documentclass[11pt,a4paper]{article}
% Kompilacja: pdfLaTeX (MiKTeX / TeX Live)
\usepackage[T1]{fontenc}
\usepackage[utf8]{inputenc}
\usepackage[polish]{babel}
\usepackage{lmodern}
\usepackage{amsmath,amssymb,amsthm,mathtools}
\usepackage{microtype}
\usepackage{xcolor}
\usepackage{geometry}
\usepackage{enumitem}
\usepackage{fancyhdr}
\usepackage[unicode]{hyperref}
\hypersetup{colorlinks=true,linkcolor=blue!45!black,urlcolor=blue!45!black}
\geometry{margin=2.5cm}
\setlength{\parindent}{0pt}
\setlength{\parskip}{4pt}
\setlength{\emergencystretch}{2em}
\setlist[enumerate]{topsep=4pt,itemsep=2pt}
\setlist[itemize]{topsep=4pt,itemsep=2pt}
\numberwithin{equation}{section}
\newcommand{\N}{\mathbb N}
\newcommand{\R}{\mathbb R}
\newcommand{\PA}{\mathrm{PA}}
\newcommand{\Q}{\mathrm Q}
\newcommand{\Prov}{\operatorname{Prov}}
\newcommand{\Prf}{\operatorname{Prf}}
\newcommand{\Bew}{\operatorname{Bew}}
\newcommand{\dotminus}{\mathbin{\dot-}}
\theoremstyle{plain}
\newtheorem{twierdzenie}{Twierdzenie}[section]
\newtheorem{lemat}[twierdzenie]{Lemat}
\newtheorem{wniosek}[twierdzenie]{Wniosek}
\newtheorem{stwierdzenie}[twierdzenie]{Stwierdzenie}
\theoremstyle{definition}
\newtheorem{definicja}[twierdzenie]{Definicja}
\newtheorem{przyklad}[twierdzenie]{Przykład}
\theoremstyle{remark}
\newtheorem{uwaga}[twierdzenie]{Uwaga}
\renewcommand{\proofname}{Dowód}
\pagestyle{fancy}
\fancyhf{}
\fancyhead[L]{\small Teoria modeli}
\fancyhead[R]{\small\ifnum\value{section}>0 Sekcja \thesection\else Wprowadzenie\fi}
\fancyfoot[C]{\thepage}
\renewcommand{\headrulewidth}{0.4pt}
\title{Teoria Modeli}
\author{}
\date{22 października 2026}
\begin{document}
\maketitle
\begin{abstract}
Niniejszy dokument przedstawia szczegółowe dowody obu twierdzeń Gödla o niezupełności.
Zaczynamy od precyzyjnych definicji systemu formalnego, arytmetyki Peana, rekurencyjności
oraz reprezentowalności. Następnie budujemy narzędzia techniczne (numerację Gödla,
lemat o reprezentacji, lemat o punkcie stałym / diagonalizacji) i na ich podstawie dowodzimy
pierwszego oraz drugiego twierdzenia o niezupełności. Reprezentowalność uzasadniamy
przez kodowanie skończonych obliczeń, a warunki wyprowadzalności --- przez wewnętrzne
weryfikowanie dowodów. Rozróżniamy wersję Gödla z 1931 r., wzmocnienie Rossera oraz
późniejsze ujęcie Hilberta--Bernaysa--Löba. Końcowe sekcje zawierają oba kierunki
twierdzenia Löwenheima--Skolema, z dowodem zwartości, oraz omówienie konsekwencji
całego skryptu: niestandardowych modeli, paradoksu Skolema, granic metod algorytmicznych
i różnicy między prawdą, dowodem oraz zaufaniem do aksjomatów.
\end{abstract}
\tableofcontents
\clearpage

% Strony 3--6 oryginalnego skanu -- transkrypcja sprawdzona na obrazach.
\input{../rozdzialy/teoria_modeli/00-wstep-i-zalozenia.tex}
\input{../rozdzialy/teoria_modeli/01-podstawy-od-alfabetu-do-teorii.tex}
\input{../rozdzialy/teoria_modeli/02-jezyk-i-system-formalny.tex}
\input{../rozdzialy/teoria_modeli/03-arytmetyka-peana-i-robinsona.tex}
\input{../rozdzialy/teoria_modeli/04-funkcje-i-zbiory-rekurencyjne.tex}
\input{../rozdzialy/teoria_modeli/05-numeracja-godla.tex}
\input{../rozdzialy/teoria_modeli/06-reprezentowalnosc.tex}
\input{../rozdzialy/teoria_modeli/07-lemat-o-punkcie-stalym-diagonalizacja.tex}
\input{../rozdzialy/teoria_modeli/08-pierwsze-twierdzenie-o-niezupelnosci.tex}
\input{../rozdzialy/teoria_modeli/09-twierdzenie-o-niedefiniowalnosci-prawdy-tarski.tex}
\input{../rozdzialy/teoria_modeli/10-warunki-hilberta-bernaysa-loba.tex}
\input{../rozdzialy/teoria_modeli/11-drugie-twierdzenie-o-niezupelnosci.tex}
\input{../rozdzialy/teoria_modeli/12-twierdzenie-loba.tex}
\input{../rozdzialy/teoria_modeli/13-twierdzenie-lowenheima-skolema.tex}
\input{../rozdzialy/teoria_modeli/14-konsekwencje-dla-modeli-niestandardowe-liczby-i-paradoks-skole.tex}
\input{../rozdzialy/teoria_modeli/15-co-niezupelnosc-oznacza-dla-codziennej-matematyki.tex}
\input{../rozdzialy/teoria_modeli/16-wlasna-niesprzecznosc-i-granice-zaufania-godel-ii-oraz-lob.tex}
\input{../rozdzialy/teoria_modeli/17-kodowanie-diagonalizacja-i-prawda-sens-twierdzenia-tarskiego.tex}
\input{../rozdzialy/teoria_modeli/18-czego-te-twierdzenia-nie-mowia.tex}
\end{document}
